\chapter{Request for Quote and Dealer Markets in Theory}\label{ch:mx:request-for-quote-and-dealer-markets-in-theory}

A client who asks five dealers for a price on a corporate bond gets a better price than one who asks two, up to a point: every dealer asked learns
that someone is selling, and the ones who lose the request can trade on it. This chapter gives the theory behind the markets that have no central
book: a search-and-bargaining model of dealer spreads, the request for quote as an auction with private values, common values and leakage, the
number of dealers a client should ask, and the networks and relationships through which dealers trade.

\section{Search and bargaining}

\begin{definition}[Over-the-counter search model, search friction, bargaining power]\label{def:mx:request-for-quote-and-dealer-markets-in-theory:search}
An \emph{over-the-counter search model}\index{over-the-counter search model} describes a market in which traders find counterparties one at a time,
at random times, and agree prices bilaterally. A \emph{search friction}\index{search friction} is the delay and cost of finding a counterparty; a
trader's \emph{bargaining power}\index{bargaining power} is the share of the gains from a trade it captures when a counterparty is found.
\end{definition}

Duffie, Gârleanu and Pedersen (2005) computed the prices at which investors trade with dealers in a dynamic model with search and bargaining, and found
that bid--ask spreads are lower when investors can more easily find other investors or reach several dealers. This chapter solves a version of their
dealer market (\texttt{firm\_otcsearch.dgp}). An asset pays one a year; investors are of a high type or a low type, the low type paying a holding cost
$\delta$; types switch at rates $\lambda_u$ (to high) and $\lambda_d$ (to low). Low-type owners sell to dealers and high-type non-owners buy from them,
meeting a dealer at rate $\rho$ (a Poisson process, One Quant Book 4, chapter~6); dealers lay off every trade in an inter-dealer market and keep a
share $z$ of each trade's surplus. With a supply $s$ below the high types' share $\lambda_u/(\lambda_u+\lambda_d)$, buyers are the long side and the
inter-dealer price is their reservation value.

In steady state the low-type owners flow in from switching and out to dealers and to the high type: $\lambda_d\mu_{ho}=(\lambda_u+\rho)\mu_{lo}$,
with $\mu_{ho}+\mu_{lo}=s$. The reservation values $\Delta V_\sigma$ (the value of owning over not owning, for type $\sigma$) solve two linear
equations (the Bellman equations of a continuous-time Markov chain, One Quant Book 4, chapter~8):
\[
r\Delta V_h=1+\lambda_d(\Delta V_l-\Delta V_h),\qquad r\Delta V_l=1-\delta+\lambda_u(\Delta V_h-\Delta V_l)+\rho(1-z)(\Delta V_h-\Delta V_l),
\]
and the ask is $\Delta V_h$, the bid $z\Delta V_l+(1-z)\Delta V_h$, the spread $z(\Delta V_h-\Delta V_l)$.

\omcode{code/firm/otcsearch/firm_otcsearch.py}{43}{58}{The dealer market's steady state and prices: the masses of owners and non-owners of each type, then the two reservation values from their linear Bellman equations, and the bid and ask that bargaining gives.}\label{lst:mx:request-for-quote-and-dealer-markets-in-theory:dgp}

With $r=5\%$, $\lambda_u=2$ and $\lambda_d=0.2$ a year, $s=0.8$, $\delta=2$ and $z=0.8$, the spread falls from 210 basis points of the price when
investors meet a dealer ten times a year to 36.7 at a hundred and 4.0 at a thousand (\cref{fig:mx:request-for-quote-and-dealer-markets-in-theory:dgp}):
faster search leaves low types less time to bear their holding cost, which is the surplus dealers bargain over.

\begin{omfigure}
\begin{tikzpicture}
\begin{axis}[width=9cm,height=5.4cm,xmode=log,ymode=log,xlabel={\small meetings with dealers a year ($\rho$)},ylabel={\small spread (bp of the price)},
  tick label style={font=\footnotesize},grid=major,grid style={gray!25},table/col sep=comma]
\addplot[omDef,very thick,mark=*] table[x=rho,y=spread_bp]{figdata/microstructure/23-request-for-quote-and-dealer-markets-in-theory/dgp.csv};
\end{axis}
\end{tikzpicture}

\omcaption{The dealer market's bid--ask spread against the investors' contact rate with dealers, in the search-and-bargaining model with the
chapter's parameters. Data: \texttt{mx\_otc.dgp\_study}.}\label{fig:mx:request-for-quote-and-dealer-markets-in-theory:dgp}
\end{omfigure}

\section{Dealer networks}

\begin{definition}[Core--periphery network]\label{def:mx:request-for-quote-and-dealer-markets-in-theory:cp}
A \emph{core--periphery network}\index{core--periphery network} of dealers has a small core of dealers who trade with each other and with many others,
and a large periphery whose dealers trade mostly with the core.
\end{definition}

Li and Schürhoff (2019) found that municipal bond dealers form such a network, with 10 to 30 hubs and over 2\,000 peripheral broker-dealers; bonds flow
from the periphery to the core and partly back; central dealers charge investors up to double the round-trip markups of peripheral ones but provide
immediacy by matching buyers and sellers more directly. The chapter's toy network (\texttt{core\_periphery}) has 10 core dealers and 200 peripheral
ones linked to two core dealers each; a bond passes along the shortest path from the seller's dealer to the buyer's, and each dealer on the path
charges a markup, 4 basis points in the core and 2 in the periphery. A client of a core dealer is reached in 1.77 hops on average and pays 9.2 basis
points; a client of a peripheral dealer, 2.58 hops and 10.4. The chain's length is part of the price; the empirical finding that the core charges more
comes from its markups, which this toy sets by hand.

\section{The request for quote as an auction}

A request for quote (One Quant Book 2, chapter~22) is a first-price sealed-bid auction among the dealers asked (One Quant Book 4, chapter~29). Each
dealer's value for the bond has a private part (its inventory, its axes: standard deviation $\sigma_p$) and it estimates the common value with an
error ($\sigma_c$). The client sells to the best bid, and each losing dealer, who now knows that a seller is about, costs the client a leakage $\ell$
on average. A dealer who wins when its estimate of the common value is the highest has overestimated it: it shades its bid by $\sigma_c\E[\max\text{ of
}n]$ (the winner's curse, One Quant Book 2, chapter~30, and chapter~22's \omterm{def:mx:portfolio-trades-and-risk-bids:agency}{risk bids}). The client's expected cost below the common value is
\[
C(n)=m-\bigl(\sqrt{\sigma_c^2+\sigma_p^2}-\sigma_c\bigr)\E[\max\text{ of }n]+\ell(n-1),
\]
with $m$ the dealers' markup: competition helps through the private values, the common part mostly cancels against the shading, and every extra dealer
adds leakage (\texttt{rfq\_cost}, on Book~2's \texttt{firm.rfq}; a simulation of the auction agrees to 0.001 basis points).

\omcode{code/firm/otcsearch/firm_otcsearch.py}{61}{70}{The client's expected cost of a request for quote to n dealers, built on Book 2's private-value formula, and the number of dealers that minimises it.}\label{lst:mx:request-for-quote-and-dealer-markets-in-theory:rfq}

\section{How many dealers to ask}

With a markup of 10 basis points, $\sigma_p=5$ and $\sigma_c=2$, \cref{fig:mx:request-for-quote-and-dealer-markets-in-theory:rfq} draws $C(n)$ for
four leakage costs, and \cref{tab:mx:request-for-quote-and-dealer-markets-in-theory:bestn} gives the optimum.

\begin{omfigure}
\begin{tikzpicture}
\begin{axis}[width=10cm,height=5.8cm,xlabel={\small dealers asked ($n$)},ylabel={\small expected cost (bp)},
  tick label style={font=\footnotesize},grid=major,grid style={gray!25},xmin=0.5,xmax=12.5,ymin=5,ymax=16,
  legend columns=4,legend style={font=\footnotesize,at={(0.5,-0.3)},anchor=north,draw=none,fill=none,
  /tikz/every even column/.append style={column sep=6pt}},table/col sep=comma]
\addplot[omDef,thick,mark=*,mark size=1.3pt] table[x=n,y=l025]{figdata/microstructure/23-request-for-quote-and-dealer-markets-in-theory/rfq.csv};
\addplot[omThm,thick,mark=square*,mark size=1.3pt] table[x=n,y=l05]{figdata/microstructure/23-request-for-quote-and-dealer-markets-in-theory/rfq.csv};
\addplot[omProp,thick,mark=triangle*,mark size=1.6pt] table[x=n,y=l1]{figdata/microstructure/23-request-for-quote-and-dealer-markets-in-theory/rfq.csv};
\addplot[gray,thick,mark=diamond*,mark size=1.6pt] table[x=n,y=l2]{figdata/microstructure/23-request-for-quote-and-dealer-markets-in-theory/rfq.csv};
\legend{{$\ell=0.25$},{$\ell=0.5$},{$\ell=1$},{$\ell=2$}}
\end{axis}
\end{tikzpicture}

\omcaption{The client's expected cost of selling by request for quote against the number of dealers asked, for four leakage costs per losing dealer
(basis points): markup 10, private-value dispersion 5, common-value error 2. Data: \texttt{mx\_otc.rfq\_study}.}\label{fig:mx:request-for-quote-and-dealer-markets-in-theory:rfq}
\end{omfigure}

\begin{table}[ht]
\centering
\small
\begin{tabular}{@{}lrrrrrr@{}}
\toprule
leakage per losing dealer (bp) & 0.1 & 0.25 & 0.5 & 1 & 2 & 3 \\
\midrule
dealers to ask & 16 & 7 & 4 & 2 & 1 & 1 \\
expected cost (bp) & 5.52 & 6.92 & 8.02 & 9.09 & 10.00 & 10.00 \\
\bottomrule
\end{tabular}
\caption{The cost-minimising number of dealers and the expected cost against the leakage cost per losing dealer. Data: \texttt{mx\_otc.rfq\_study}.}\label{tab:mx:request-for-quote-and-dealer-markets-in-theory:bestn}
\end{table}

The optimum falls fast: from 16 dealers when a loser costs a tenth of a basis point to 4 at half a basis point and a single dealer from 2 basis
points. The common-value part matters too: at half a basis point of leakage, a client should ask 6 dealers if the dealers' estimates had no common
error, 4 with $\sigma_c=2$ and 3 with $\sigma_c=5$, because the common part adds noise that the winner shades away without adding competition.
Riggs, Onur, Reiffen and Zhu (2020) found that customers on the two largest dealer-to-customer swap execution facilities for index CDS typically contact
few dealers, and modelled why: the benefit of wider exposure is mitigated by the winner's curse and by customer--dealer relationships.

\section{Relationships and information chasing}

\begin{definition}[Information chasing]\label{def:mx:request-for-quote-and-dealer-markets-in-theory:chasing}
\emph{Information chasing}\index{information chasing} is a dealer's quoting of better prices to clients whose trades carry information, because
what it learns from their flow is worth more to it, in its trading with others, than the adverse selection it suffers on their trades.
\end{definition}

In a competitive quote, a dealer's spread to a client is its cost plus the adverse selection it expects, less the value of what it learns. With a
cost of 2 basis points and adverse selection of 1.5 on an informed client's trade, the informed client is quoted 3.5 when its information is worth
nothing to the dealer, the same 2.0 as an uninformed client when it is worth 1.5, and 0.5 when it is worth 3 (\texttt{chase}): the informed client
gets the better price. Relationships work the same way: Hendershott, Li, Livdan and Schürhoff (2020) found in insurers' corporate bond trades that
execution prices first improve with the number of dealers in a client's network and worsen beyond about twenty, as repeat business with a few
dealers trades off against competition among many; Di Maggio, Kermani and Song (2017) found that dealers charge lower spreads to the dealers with whom
they have the strongest ties, more so in turmoil, and that systemically important dealers charge peripheral dealers and clients more.

\section{Tutorial: dealers, searches and auctions}\label{tut:mx:request-for-quote-and-dealer-markets-in-theory:tut}

\begin{tutorial}
\textbf{Goal.} Solve a search-and-bargaining dealer market, find the cost-minimising number of dealers in a request for quote, and trace
intermediation chains in a \omterm{def:mx:request-for-quote-and-dealer-markets-in-theory:cp}{core--periphery network}.
\textbf{End state:} \cref{fig:mx:request-for-quote-and-dealer-markets-in-theory:dgp,fig:mx:request-for-quote-and-dealer-markets-in-theory:rfq},
\cref{tab:mx:request-for-quote-and-dealer-markets-in-theory:bestn} and the numbers of sections 2 and 5.
\begin{tutsteps}
\item \textbf{Search.} \texttt{firm\_otcsearch.dgp(lam\_u, lam\_d, rho, s, r, delta, z)}; \texttt{mx\_otc.dgp\_study()}.
\item \textbf{Auction.} \texttt{rfq\_cost}, \texttt{best\_n}, \texttt{rfq\_simulate}; \texttt{rfq\_study()}, \texttt{common\_study()}.
\item \textbf{Networks and chasing.} \texttt{core\_periphery}, \texttt{intermediation}, \texttt{chase}; \texttt{network\_study()},
\texttt{chase\_study()}; draw with \texttt{fig\_otc.py}.
\end{tutsteps}
\textbf{What to change next.} Let investors also meet each other directly (Duffie, Gârleanu and Pedersen's other channel); make leakage depend on
the losers' inventories; let peripheral dealers choose their core links.
\end{tutorial}

\section{Build: OTC search and auctions}\label{bld:mx:request-for-quote-and-dealer-markets-in-theory:otcsearch}

\begin{build}
\bfield{Purpose} The theory behind chapter~21's bond adapter and the desk's dealer choices: prices from search, the right number of dealers, and
intermediation chains.
\bfield{Interface} \texttt{dgp(lam\_u, lam\_d, rho, s, r, delta, z)}, \texttt{rfq\_cost(n, markup, sigma\_p, sigma\_c, leak)}, \texttt{best\_n},
\texttt{rfq\_simulate}, \texttt{chase(cost, adverse, info\_value)}, \texttt{core\_periphery(n\_core, n\_periph, links, seed)},
\texttt{intermediation(adj, n\_core, markup\_core, markup\_periph, trials, seed)}.
\bfield{Rules} The asset pays one a year; costs in basis points below the common value; the version of the search model with buyers on the long side.
\bfield{Acceptance tests} \texttt{code/firm/otcsearch/tests/}: the steady state's flow balance and masses; the spread falls with the contact rate and
vanishes without dealer \omterm{def:mx:request-for-quote-and-dealer-markets-in-theory:search}{bargaining power}; the RFQ formula against Book~2's and against simulation; the optimum's monotonicity in leakage; chasing
lowers the informed spread; peripheral trades take longer chains.
\bfield{Stretch} Investor-to-investor search; endogenous dealer entry; inventory-dependent leakage.
\end{build}

\begin{omsources}
\item D.~Duffie, N.~Gârleanu and L.~H.~Pedersen, ``Over-the-counter markets'', \emph{Econometrica} 73(6), 2005.
\item P.~Di~Maggio, A.~Kermani and Z.~Song, ``The value of trading relations in turbulent times'', \emph{Journal of Financial Economics} 124(2),
2017.
\item D.~Li and N.~Schürhoff, ``Dealer networks'', \emph{Journal of Finance} 74(1), 2019.
\item L.~Riggs, E.~Onur, D.~Reiffen and H.~Zhu, ``Swap trading after Dodd--Frank: evidence from index CDS'', \emph{Journal of Financial Economics}
137(3), 2020.
\item T.~Hendershott, D.~Li, D.~Livdan and N.~Schürhoff, ``Relationship trading in over-the-counter markets'', \emph{Journal of Finance} 75(2), 2020.
\end{omsources}

\section{Exercises}

\begin{exercise}[$\star$]\label{exo:mx:request-for-quote-and-dealer-markets-in-theory:1}
With $\lambda_u=2$, $\lambda_d=0.2$, $s=0.8$ and $\rho=100$, compute the mass of low-type owners in steady state.
\end{exercise}

\begin{exercise}[$\star$]\label{exo:mx:request-for-quote-and-dealer-markets-in-theory:2}
Compute $C(n)$ for $n=1$ and $n=2$ with the chapter's parameters and a leakage of 1 basis point.
\end{exercise}

\begin{exercise}[$\star$]\label{exo:mx:request-for-quote-and-dealer-markets-in-theory:3}
At what value of an informed trade's information does a dealer quote informed and uninformed clients the same spread?
\end{exercise}

\begin{exercise}[$\star\star$]\label{exo:mx:request-for-quote-and-dealer-markets-in-theory:4}
Why does the spread vanish when dealers have no \omterm{def:mx:request-for-quote-and-dealer-markets-in-theory:search}{bargaining power}, and why does it fall with the contact rate?
\end{exercise}

\begin{exercise}[$\star\star$]\label{exo:mx:request-for-quote-and-dealer-markets-in-theory:5}
Why does a common-value error reduce the number of dealers a client should ask?
\end{exercise}

\begin{exercise}[$\star\star$]\label{exo:mx:request-for-quote-and-dealer-markets-in-theory:6}
In the toy network, why do clients of peripheral dealers go through longer chains, and what would make them pay less?
\end{exercise}

\begin{exercise}[$\star\star\star$]\label{exo:mx:request-for-quote-and-dealer-markets-in-theory:7}
\emph{Coding.} With no common-value error and a leakage of 1 basis point, how many dealers should the client ask, and what does it cost?
\end{exercise}

\begin{exercise}[$\star\star\star$]\label{exo:mx:request-for-quote-and-dealer-markets-in-theory:8}
\emph{Find the flaw.} ``More dealers always means more competition, so our RFQs go to every dealer on the platform.''
\end{exercise}

\section{Problem: How Many Dealers to Ask}

\begin{problem}[{Weekend problem --- how many dealers to ask}]\label{pb:mx:request-for-quote-and-dealer-markets-in-theory:1}
A fund's bond desk sends every RFQ to eight dealers. The head of trading asks whether that is right.

\medskip\noindent\textbf{Part I --- Search.}
\begin{enumerate}
\item Define a \omterm{def:mx:request-for-quote-and-dealer-markets-in-theory:search}{search friction} and \omterm{def:mx:request-for-quote-and-dealer-markets-in-theory:search}{bargaining power}.
\item State the search model's types, rates and trades.
\item Derive the steady-state mass of low-type owners.
\item Write the Bellman equations for the reservation values and the bid and ask.
\item How does the spread depend on the contact rate?
\end{enumerate}

\medskip\noindent\textbf{Part II --- The auction.}
\begin{enumerate}[resume]
\item Why is an RFQ a first-price sealed-bid auction?
\item Derive the winner's-curse shading for a common-value error.
\item Write the client's expected cost $C(n)$ and explain each term.
\item How closely does a simulation of the auction agree?
\end{enumerate}

\medskip\noindent\textbf{Part III --- The number of dealers.}
\begin{enumerate}[resume]
\item \emph{State the named result}: the cost-minimising number of dealers and its dependence on the leakage cost per losing dealer.
\item How does the common-value error change it?
\item What did Riggs, Onur, Reiffen and Zhu find about how many dealers customers contact, and why?
\item Is eight dealers right for this desk? What would you need to know?
\end{enumerate}

\medskip\noindent\textbf{Part IV --- Networks and relationships.}
\begin{enumerate}[resume]
\item Define a \omterm{def:mx:request-for-quote-and-dealer-markets-in-theory:cp}{core--periphery network} and give Li and Schürhoff's findings.
\item What does the toy network show about chains?
\item Define \omterm{def:mx:request-for-quote-and-dealer-markets-in-theory:chasing}{information chasing} and give the spreads it implies.
\item What did Hendershott, Li, Livdan and Schürhoff find about network size and prices?
\item What did Di Maggio, Kermani and Song find about dealer ties in turmoil?
\item How would you measure the leakage per losing dealer from a desk's own data?
\item In one sentence: what does a client trade off when it asks one more dealer?
\end{enumerate}
\end{problem}

\section{Interview questions}

\begin{interviewq}[$\star$ \iqroles{trader}]\label{iq:mx:request-for-quote-and-dealer-markets-in-theory:1}
You need to sell an illiquid corporate bond. How many dealers do you ask, and why?
\end{interviewq}

\begin{interviewq}[$\star\star$ \iqroles{researcher}]\label{iq:mx:request-for-quote-and-dealer-markets-in-theory:2}
Explain the winner's curse in an RFQ and how a dealer should bid.
\end{interviewq}

\begin{interviewq}[$\star\star$ \iqroles{researcher}]\label{iq:mx:request-for-quote-and-dealer-markets-in-theory:3}
In a search model, why are spreads lower when investors can find dealers more easily?
\end{interviewq}

\begin{interviewq}[$\star\star$ \iqroles{bank}]\label{iq:mx:request-for-quote-and-dealer-markets-in-theory:4}
Why might a dealer quote a tighter price to a client it knows is informed?
\end{interviewq}

\begin{interviewq}[$\star\star$ \iqroles{developer}]\label{iq:mx:request-for-quote-and-dealer-markets-in-theory:5}
Design the data an RFQ platform should keep to let clients measure leakage.
\end{interviewq}

\begin{interviewq}[$\star\star\star$ \iqroles{risk}]\label{iq:mx:request-for-quote-and-dealer-markets-in-theory:6}
The core dealers of a bond market are hit by a shock. What happens to intermediation chains and to clients' costs?
\end{interviewq}
